%============================================================
%  Bd. 17 - Folgen
%============================================================

\documentclass[main.tex]{subfiles}

\ifSubfilesClassLoaded{
  \BandSetupStandalone{B17}
}{
}

\title{Bd. 17 - Folgen}
\author{Martin Kunze}
\date{}
\setcounter{file}{17}

\hypersetup{
  pdftitle={Bd. 17 - Folgen},
  pdfauthor={Martin Kunze}
}

\providecommand{\PosNatSeg}[1]{\mathbb{N}_{1,\ldots,#1}}
% Die folgenden Symbole werden in Band 15 lokal eingeführt. Die Fallbacks
% halten diesen Band auch im Standalone-Build lesbar und kollidieren wegen
% \providecommand nicht mit den dort bereits vorhandenen Definitionen.
\ifSubfilesClassLoaded{%
  \providecommand{\RealLe}{\mathrel{\leq_{\mathbb R}}}
  \providecommand{\RealLt}{\mathrel{<_{\mathbb R}}}
  \providecommand{\RealZero}{0_{\mathbb R}}
  \providecommand{\RealOne}{1_{\mathbb R}}
  \providecommand{\RealAdd}{\mathbin{\oplus}}
  \providecommand{\RealMul}{\mathbin{\odot}}
  \providecommand{\RealNeg}[1]{\mathop{\ominus}\!#1}
  \providecommand{\RealInv}[1]{\operatorname{inv}_{\mathbb R}\!\left(#1\right)}
  \providecommand{\RealAbs}[1]{\left\lvert #1\right\rvert_{\mathbb R}}
  \providecommand{\RealDist}[2]{d_{\mathbb R}\!\left(#1,#2\right)}
}{%
}

\begin{document}

\title{Bd. 17 - Folgen}
\author{Martin Kunze}
\date{}
\hypersetup{pageanchor=false}
\maketitle
\hypersetup{pageanchor=true}
\setcounter{tocdepth}{3}
\tableofcontents
\setcounter{file}{17}
\renewcommand{\FormulaBandID}{17}

\chapter{Einleitung}

Eine Folge ist keine neue Art mathematischer Gegenstand, sondern eine
Funktion, deren Definitionsmenge als Indexmenge gelesen wird. Diese Sichtweise
trennt drei Dinge sauber voneinander: die Indexmenge, die Zielmenge und die
tatsächlich angenommene Bildmenge. Insbesondere darf aus einer Schreibweise
\(a\colon I\to M\) nicht ohne einen Surjektivitätsbeweis geschlossen werden,
dass jedes Element von \(M\) als Folgenglied vorkommt.

Der Band beginnt deshalb mit beliebig indizierten Familien. Unendliche Folgen
\(\mathbb N\to M\), endliche Folgen und die traditionelle Schreibweise
\((a_i)_{i\in I}\) erscheinen anschließend als Spezialfälle. Für reelle
Folgen werden Konvergenz, Grenzwerte, Beschränktheit, die elementaren
Grenzwertregeln und das Cauchy-Kriterium entwickelt. Das letzte Kapitel stellt
die rein folgen- und mengentheoretischen Familien bereit, die in der
Mogiljanskaja-Konstruktion der späteren Halbgruppentheorie verwendet werden.

\chapter{Indizierte Familien}

\section{Familien als Funktionen}

\FormulaDefDeltaK[Indizierte Familie]{%
  a\colon I\to M\vdash
  (a_i)_{i\in I}\coloneqq a\dsep
  i\in I\vdash a_i\coloneqq a(i)%
}{IndexedFamilyDef}{%
  \DeltaRow{Mengen}{I\dsep M\dsep i}
  \DeltaRow{Funktionen}{a}
  \DeltaRow{Neue Schreibweisen}{(a_i)_{i\in I}\dsep a_i}
}

Die Menge \(I\) heißt \textbf{Indexmenge}, \(M\) heißt \textbf{Zielmenge}
der Familie, und \(a_i\) heißt das zum Index \(i\) gehörende
\textbf{Familienglied}. Die Klammer \((a_i)_{i\in I}\) bezeichnet die ganze
Funktion und nicht die ungeordnete Menge ihrer Werte. Verschiedene Indizes
dürfen daher dasselbe Familienglied tragen.

Im folgenden Metaschema bezeichnet \(t(i)\) kein zusätzliches freies
Funktionssymbol, sondern eine bereits wohldefinierte, durch \(i\) eindeutig
bestimmte Funktionsvorschrift. Dies ist genau die Vorschriftsprämisse des
Funktionskonstruktionssatzes aus Band~05; bei einem gewöhnlichen Term bleibt
als inhaltlicher Nachweis daher nur noch seine Zugehörigkeit zur Zielmenge.

\FormulaThmDeltaK[Einmaliges Definitions- und Typisierungsprinzip für
Familien]{%
  \forall i\in I\,\bigl(t(i)\in M\bigr)
  \vdash
  \exists!a\colon I\to M\,
    \forall i\in I\,\bigl(a_i=t(i)\bigr)%
}{IndexedFamilyTypedDefinitionPrinciple}{%
  \DeltaRow{Mengen}{I\dsep M\dsep i}
  \DeltaRow{Funktionen}{a}
  \DeltaRow{Termschema}{t}
  \DeltaRow{Metavoraussetzung}{t(i)\text{ ist für jedes }i
    \text{ eindeutig bestimmt}}
}
\begin{proof}
Die Metavoraussetzung liefert die eindeutige Funktionsvorschrift, und die
angezeigte Voraussetzung zeigt, dass ihr Wert in \(M\) liegt. Das allgemeine
Konstruktionsprinzip für Funktionen aus Band~05 liefert daher
genau eine Funktion \(a\colon I\to M\) mit \(a(i)=t(i)\). Mit
\FormulaRefAuto{IndexedFamilyDef}[def] ist dies gleichbedeutend mit
\(a_i=t(i)\) für alle \(i\in I\).
\end{proof}

Dieses Prinzip ist die maßgebliche Typisierungsregel für Folgen: Bei der
Definition wird einmal für einen beliebigen Index \(i\in I\) nachgewiesen,
dass der vorgeschriebene Wert in \(M\) liegt. Danach ist die Typisierung in
der Funktionsaussage \(a\colon I\to M\) gespeichert und muss nicht bei jeder
Verwendung eines Gliedes wiederholt werden.

\FormulaThmDeltaK[Automatische Typisierung eines Familiengliedes]{%
  a\colon I\to M\dsep i\in I\vdash a_i\in M%
}{IndexedFamilyCoordinateTyping}{%
  \DeltaRow{Mengen}{I\dsep M\dsep i}
  \DeltaRow{Funktionen}{a}
}
\begin{proof}
Aus der Typisierung der Funktion folgt \(a(i)\in M\). Nach der
Koordinatendefinition ist \(a_i=a(i)\).
\end{proof}

\FormulaThmDeltaK[Erweiterung der Zielmenge einer Familie]{%
  a\colon I\to M\dsep M\subseteq N\vdash a\colon I\to N%
}{IndexedFamilyCodomainExtension}{%
  \DeltaRow{Mengen}{I\dsep M\dsep N}
  \DeltaRow{Funktionen}{a}
}
\begin{proof}
Dies ist die Zielmengenerweiterung aus Band~05, auf die als Familie gelesene
Funktion \(a\) angewandt. Die Werte und damit alle Glieder \(a_i\) ändern
sich dabei nicht.
\end{proof}

\section{Bildmenge und Wertemenge}

\FormulaDefDeltaK[Wertemenge einer Familie]{%
  a\colon I\to M\vdash
  \operatorname{Bild}(a)\coloneqq a[I]%
}{IndexedFamilyImageDef}{%
  \DeltaRow{Mengen}{I\dsep M\dsep a[I]}
  \DeltaRow{Funktionen}{a}
  \DeltaRow{Neue Schreibweise}{\operatorname{Bild}(a)}
}

\FormulaThmDeltaK[Elementkriterium der Wertemenge]{%
  a\colon I\to M\vdash
  x\in\operatorname{Bild}(a)
  \leftrightarrow
  \exists i\in I\,(x=a_i)%
}{IndexedFamilyImageMembership}{%
  \DeltaRow{Mengen}{I\dsep M\dsep x\dsep i}
  \DeltaRow{Funktionen}{a}
}
\begin{proof}
Nach der Definition der Bildmenge einer Teilmenge in Band~05 gilt
\(x\in a[I]\) genau dann, wenn ein \(i\in I\) mit \(x=a(i)\) existiert.
Nun wird \(a(i)\) durch \(a_i\) ersetzt.
\end{proof}

\FormulaThmDeltaK[Wertemenge liegt in der Zielmenge]{%
  a\colon I\to M\vdash \operatorname{Bild}(a)\subseteq M%
}{IndexedFamilyImageSubsetCodomain}{%
  \DeltaRow{Mengen}{I\dsep M}
  \DeltaRow{Funktionen}{a}
}

\FormulaThmDeltaK[Surjektive Familien und ihre Wertemenge]{%
  a\colon I\to M\vdash
  \bigl(a\colon I\sur M\bigr)
  \leftrightarrow
  \operatorname{Bild}(a)=M%
}{IndexedFamilySurjectiveIffImage}{%
  \DeltaRow{Mengen}{I\dsep M}
  \DeltaRow{Funktionen}{a}
}

Die Zielmenge gehört zur Typisierung, die Wertemenge wird dagegen durch die
Glieder bestimmt. So kann dieselbe Familie zunächst als \(a\colon I\to M\)
und nach einer Zielmengenerweiterung als \(a\colon I\to N\) gelesen werden,
ohne dass sich \(\operatorname{Bild}(a)\) ändert.

\section{Gleichheit, Einschränkung und Umindizierung}

\FormulaThmDeltaK[Extensionalität indizierter Familien]{%
  a,b\colon I\to M\vdash
  a=b\leftrightarrow\forall i\in I\,(a_i=b_i)%
}{IndexedFamilyExtensionality}{%
  \DeltaRow{Mengen}{I\dsep M\dsep i}
  \DeltaRow{Funktionen}{a\dsep b}
}
\begin{proof}
Dies ist die Funktionsextensionalität aus Band~05 in Koordinatenschreibweise.
Insbesondere genügt die Gleichheit der Wertemengen nicht für die Gleichheit
zweier Familien.
\end{proof}

\FormulaDefDeltaK[Einschränkung einer Familie]{%
  a\colon I\to M\dsep J\subseteq I\vdash
  (a_i)_{i\in J}\coloneqq a\restriction_J%
}{IndexedFamilyRestrictionDef}{%
  \DeltaRow{Mengen}{I\dsep J\dsep M\dsep i}
  \DeltaRow{Funktionen}{a\dsep a\restriction_J}
}

\FormulaThmDeltaK[Eigenschaften der eingeschränkten Familie]{%
  \begin{aligned}[t]
  &a\colon I\to M\dsep J\subseteq I\vdash
    a\restriction_J\colon J\to M\dsep{}\\[-2pt]
  &\quad\forall i\in J\,\bigl((a\restriction_J)_i=a_i\bigr)\dsep{}\\[-2pt]
  &\quad\operatorname{Bild}(a\restriction_J)\subseteq
    \operatorname{Bild}(a)
  \end{aligned}%
}{IndexedFamilyRestrictionFacts}{%
  \DeltaRow{Mengen}{I\dsep J\dsep M\dsep i}
  \DeltaRow{Funktionen}{a\dsep a\restriction_J}
}

\FormulaDefDeltaK[Umindizierte Familie]{%
  a\colon I\to M\dsep \varphi\colon J\to I\vdash
  (a_{\varphi(j)})_{j\in J}\coloneqq a\circ\varphi%
}{IndexedFamilyReindexingDef}{%
  \DeltaRow{Mengen}{I\dsep J\dsep M\dsep j}
  \DeltaRow{Funktionen}{a\dsep\varphi\dsep a\circ\varphi}
}

\FormulaThmDeltaK[Bildmenge nach Umindizierung]{%
  a\colon I\to M\dsep \varphi\colon J\to I\vdash
  \operatorname{Bild}(a\circ\varphi)=a[\varphi[J]]
  \subseteq\operatorname{Bild}(a)%
}{IndexedFamilyReindexingImage}{%
  \DeltaRow{Mengen}{I\dsep J\dsep M}
  \DeltaRow{Funktionen}{a\dsep\varphi}
}

Ist \(\varphi\) injektiv, so heißt die umindizierte Familie eine
\textbf{Teilfamilie}. Eine bloße Einschränkung ist der Spezialfall, in dem
\(\varphi\) die Inklusion \(J\hookrightarrow I\) ist. Ist \(\varphi\)
surjektiv, dann hat die umindizierte Familie dieselbe Wertemenge wie die
Ausgangsfamilie.

\section{Injektive und surjektive Familien}

\FormulaDefDeltaK[Injektive Familie]{%
  a\colon I\to M\vdash
  \bigl((a_i)_{i\in I}\text{ ist injektiv}\bigr)
  \coloneqq \bigl(a\colon I\inj M\bigr)%
}{InjectiveIndexedFamilyDef}{%
  \DeltaRow{Mengen}{I\dsep M}
  \DeltaRow{Funktionen}{a}
}

\FormulaThmDeltaK[Indexkriterium für Injektivität]{%
  a\colon I\to M\vdash
  \bigl(a\colon I\inj M\bigr)
  \leftrightarrow
  \forall i,j\in I\,(a_i=a_j\rightarrow i=j)%
}{InjectiveIndexedFamilyCriterion}{%
  \DeltaRow{Mengen}{I\dsep M\dsep i\dsep j}
  \DeltaRow{Funktionen}{a}
}

\FormulaDefDeltaK[Surjektive Familie]{%
  a\colon I\to M\vdash
  \bigl((a_i)_{i\in I}\text{ durchläuft }M\bigr)
  \coloneqq \bigl(a\colon I\sur M\bigr)%
}{SurjectiveIndexedFamilyDef}{%
  \DeltaRow{Mengen}{I\dsep M}
  \DeltaRow{Funktionen}{a}
}

Eine bijektive Familie ist zugleich injektiv und surjektiv. Sie zählt jedes
Element ihrer Zielmenge genau einmal auf. Eine nichtinjektive Folge darf
dagegen Werte wiederholen, und eine nichtsurjektive Folge lässt Zielwerte aus.

\chapter{Endliche und unendliche Folgen}

\section{Unendliche Folgen}

\FormulaDefDeltaK[Unendliche Folge in einer Menge]{%
  a\colon\mathbb N\to M\vdash
  (a_n)_{n\in\mathbb N}\coloneqq a%
}{InfiniteSequenceDef}{%
  \DeltaRow{Mengen}{M\dsep n}
  \DeltaRow{Funktionen}{a}
  \DeltaRow{Neue Schreibweise}{(a_n)_{n\in\mathbb N}}
}

Unsere natürlichen Zahlen beginnen bei \(0\). Daher ist \(a_0\) das erste,
\(a_1\) das zweite und allgemein \(a_n\) das durch den Index \(n\)
bezeichnete Glied. Die drei Schreibweisen
\[
  a\colon\mathbb N\to M,
  \qquad (a_n)_{n\in\mathbb N},
  \qquad (a_0,a_1,a_2,\ldots)
\]
bezeichnen dieselbe Funktion, betonen aber unterschiedliche Aspekte.

\FormulaDefDeltaK[Verschobene Folge und Restfolge]{%
  a\colon\mathbb N\to M\dsep r\in\mathbb N\vdash
  a^{\langle r\rangle}\colon\mathbb N\to M\dsep
  a^{\langle r\rangle}_n\coloneqq a_{n+r}%
}{SequenceTailDef}{%
  \DeltaRow{Mengen}{M\dsep n\dsep r}
  \DeltaRow{Funktionen}{a\dsep a^{\langle r\rangle}}
}

Die Restfolge lässt genau die ersten \(r\) Glieder aus. Ihre Wertemenge ist
im Allgemeinen nur eine Teilmenge der ursprünglichen Wertemenge.

\FormulaDefDeltaK[Teilfolge einer unendlichen Folge]{%
  a\colon\mathbb N\to M\dsep
  \varphi\colon\mathbb N\to\mathbb N\dsep
  \forall n\in\mathbb N\,
    \bigl(\varphi(n)<\varphi(\Succ(n))\bigr)
  \vdash
  (a_{\varphi(n)})_{n\in\mathbb N}\coloneqq a\circ\varphi%
}{SubsequenceDef}{%
  \DeltaRow{Mengen}{M\dsep n}
  \DeltaRow{Funktionen}{a\dsep\varphi\dsep a\circ\varphi}
}

Die strikte Zunahme der Indexfolge verhindert Wiederholungen und eine
Umkehrung der ursprünglichen Reihenfolge. Jede Restfolge ist eine Teilfolge;
die zugehörige Indexfunktion lautet \(\varphi(n)=n+r\).

\section{Endliche Folgen mit nullbasierten Indizes}

Für \(n\in\mathbb N\) ist
\(\NatSeg{n}=\{i\in\mathbb N\mid i<n\}\) der Anfangsabschnitt mit genau
\(n\) Elementen. Insbesondere ist \(\NatSeg{0}=\varnothing\), während
\(\NatSeg{1}=\{0\}\) gilt.

\FormulaDefDeltaK[Nullbasierte endliche Folge]{%
  n\in\mathbb N\dsep a\colon\NatSeg{n}\to M\vdash
  (a_i)_{i<n}\coloneqq(a_i)_{i\in\NatSeg{n}}%
}{ZeroBasedFiniteSequenceDef}{%
  \DeltaRow{Mengen}{M\dsep n\dsep i\dsep\NatSeg{n}}
  \DeltaRow{Funktionen}{a}
}

Eine nullbasierte Folge der Länge \(n\) wird für \(n>0\) auch als
\[
  (a_0,a_1,\ldots,a_{n-1})
\]
geschrieben. Der größte Index ist also \(n-1\), die Anzahl der Glieder aber
\(n\). Bei \(n=0\) gibt es keinen Ausdruck \(a_{n-1}\); die Folge ist dann
die eindeutig bestimmte leere Funktion \(\varnothing\to M\).

\section{Endliche Folgen mit positiven Indizes}

\FormulaDefDeltaK[Positive Indexmenge]{%
  n\in\mathbb N\vdash
  \PosNatSeg{n}\coloneqq
  \{i\in\mathbb N\mid 1\leq i\land i\leq n\}%
}{PositiveNaturalSegmentDef}{%
  \DeltaRow{Mengen}{n\dsep i\dsep\PosNatSeg{n}}
  \DeltaRow{Neue Schreibweise}{\PosNatSeg{n}}
}

\FormulaThmDeltaK[Ränder der positiven Indexmenge]{%
  n\in\mathbb N\vdash
  \PosNatSeg{0}=\varnothing\land
  \forall i\in\mathbb N\,
    \bigl(i\in\PosNatSeg{n}\leftrightarrow
      1\leq i\land i\leq n\bigr)%
}{PositiveNaturalSegmentBoundary}{%
  \DeltaRow{Mengen}{n\dsep i\dsep\PosNatSeg{n}}
}
\begin{proof}
Beide Aussagen folgen unmittelbar aus
\FormulaRefAuto{PositiveNaturalSegmentDef}[def]. Für \(n=0\) kann keine
natürliche Zahl zugleich \(1\leq i\) und \(i\leq0\) erfüllen.
\end{proof}

\FormulaDefDeltaK[Positiv indizierte endliche Folge]{%
  n\in\mathbb N\dsep a\colon\PosNatSeg{n}\to M\vdash
  (a_i)_{i=1}^{n}\coloneqq(a_i)_{i\in\PosNatSeg{n}}%
}{PositiveFiniteSequenceDef}{%
  \DeltaRow{Mengen}{M\dsep n\dsep i\dsep\PosNatSeg{n}}
  \DeltaRow{Funktionen}{a}
}

Für \(n>0\) lautet die ausgeschriebene Form
\[
  (a_1,a_2,\ldots,a_n).
\]
Auch diese Folge hat genau \(n\) Glieder. Die Schreibweise
\((a_0,\ldots,a_n)\) hätte dagegen \(n+1\) Glieder und darf nicht als Folge
der Länge \(n\) bezeichnet werden.

\FormulaDefDeltaK[Indexverschiebung zwischen den beiden Konventionen]{%
  n\in\mathbb N\vdash
  \lambda_n\colon\NatSeg{n}\to\PosNatSeg{n}\dsep
  i\in\NatSeg{n}\vdash\lambda_n(i)\coloneqq\Succ(i)%
}{FiniteSequenceIndexShiftDef}{%
  \DeltaRow{Mengen}{n\dsep i\dsep\NatSeg{n}\dsep\PosNatSeg{n}}
  \DeltaRow{Funktionen}{\lambda_n}
}
\begin{proof}
Aus \(i\in\NatSeg{n}\) folgen \(i<n\) und damit
\(1\leq\Succ(i)\leq n\), also
\(\Succ(i)\in\PosNatSeg{n}\). Das einmalige Typisierungsprinzip liefert die
angezeigte Funktion.
\end{proof}

\FormulaThmDeltaK[Äquivalenz der beiden endlichen Indexkonventionen]{%
  n\in\mathbb N\vdash
  \lambda_n\colon\NatSeg{n}\bij\PosNatSeg{n}\dsep
  j\in\PosNatSeg{n}\vdash\lambda_n^{-1}(j)=\Pred(j)%
}{FiniteSequenceIndexShiftBijective}{%
  \DeltaRow{Mengen}{n\dsep i\dsep j\dsep\NatSeg{n}\dsep\PosNatSeg{n}}
  \DeltaRow{Funktionen}{\lambda_n\dsep\lambda_n^{-1}}
}
\begin{proof}
Für \(i<n\) gilt \(1\leq\Succ(i)\leq n\). Umgekehrt besitzt jedes
\(j\) mit \(1\leq j\leq n\) den eindeutig bestimmten Vorgänger
\(\Pred(j)<n\). Nachfolger und Vorgänger sind auf diesen Bereichen invers.
\end{proof}

\FormulaThmDeltaK[Umrechnung endlicher Folgen]{%
  \begin{aligned}[t]
  &n\in\mathbb N\dsep a\colon\NatSeg{n}\to M\vdash
    a\circ\lambda_n^{-1}\colon\PosNatSeg{n}\to M\dsep{}\\[-2pt]
  &\quad\forall j\in\PosNatSeg{n}\,
    \bigl((a\circ\lambda_n^{-1})_j=a_{\Pred(j)}\bigr)
  \end{aligned}%
}{FiniteSequenceConventionConversion}{%
  \DeltaRow{Mengen}{M\dsep n\dsep j}
  \DeltaRow{Funktionen}{a\dsep\lambda_n}
}

Die Umrechnung ändert nur die Namen der Indizes, weder die Reihenfolge noch
die Wertemenge. Damit können nullbasierte und positiv indizierte endliche
Folgen verwendet werden, ohne eine versteckte Verschiebung der Länge
einzuführen.

\chapter{Grundlegende Beispiele und Rekursion}

\section{Konstante und kanonische Familien}

\FormulaDefDeltaK[Konstante Familie]{%
  c\in M\vdash
  \operatorname{const}_{I,c}\colon I\to M\dsep
  i\in I\vdash(\operatorname{const}_{I,c})_i\coloneqq c%
}{ConstantIndexedFamilyDef}{%
  \DeltaRow{Mengen}{I\dsep M\dsep c\dsep i}
  \DeltaRow{Funktionen}{\operatorname{const}_{I,c}}
}

Ist \(I\neq\varnothing\), so hat die konstante Familie die Wertemenge
\(\{c\}\); bei leerer Indexmenge ist ihre Wertemenge leer. Für
\(I=\mathbb N\) entsteht die konstante Folge \((c,c,c,\ldots)\).

\FormulaDefDeltaK[Kanonische Aufzählungsfolge der natürlichen Zahlen]{%
  \iota_{\mathbb N}\colon\mathbb N\to\mathbb N\dsep
  n\in\mathbb N\vdash(\iota_{\mathbb N})_n\coloneqq n%
}{NaturalEnumerationSequenceDef}{%
  \DeltaRow{Mengen}{n}
  \DeltaRow{Funktionen}{\iota_{\mathbb N}}
}

Diese Folge ist bijektiv und hat Bildmenge \(\mathbb N\). Die Folge
\((0,2,4,\ldots)\) der geraden natürlichen Zahlen entsteht dagegen durch
die Vorschrift \(e_n=2n\); ihre Zielmenge kann \(\mathbb N\) sein, ihre
Wertemenge ist aber nur die Menge der geraden Zahlen.

\section{Rekursiv definierte Folgen}

\FormulaThmDeltaK[Rekursive Definition einer Folge durch einen Schritt]{%
  x_0\in M\dsep F\colon M\to M
  \vdash
  \exists!a\colon\mathbb N\to M\,
  \bigl(a_0=x_0\land
    \forall n\in\mathbb N\,(a_{\Succ(n)}=F(a_n))\bigr)%
}{RecursiveSequenceExistenceUnique}{%
  \DeltaRow{Mengen}{M\dsep x_0\dsep n}
  \DeltaRow{Funktionen}{F\dsep a}
}
\begin{proof}
Dies ist der Rekursionssatz für die natürlichen Zahlen aus Band~12,
formuliert in Folgennotation. Die Voraussetzung \(F\colon M\to M\) enthält
bereits den einmaligen Nachweis, dass jeder Rekursionsschritt wieder in
\(M\) liegt.
\end{proof}

\begin{example}[Iteration]
Zu \(x_0\in M\) und \(F\colon M\to M\) gehört die eindeutig bestimmte
Iterationsfolge
\[
  x_0,\quad F(x_0),\quad F(F(x_0)),\quad\ldots.
\]
Man schreibt auch \(a_n=F^n(x_0)\), wobei \(F^0=\Id_M\) ist.
\end{example}

\begin{example}[Arithmetische Folge]
Für \(u,d\in\mathbb R\) ist nach Band~15 die Vorschrift
\(F_d(x)=x\RealAdd d\) eine Selbstabbildung von \(\mathbb R\). Daher wird
durch
\[
  a_0=u,
  \qquad a_{\Succ(n)}=F_d(a_n)
\]
eine eindeutig bestimmte reelle Folge definiert. In der konventionellen
Notation des nächsten Kapitels lautet der Schritt \(a_{n+1}=a_n+d\).
\end{example}

\begin{example}[Geometrische Folge]
Für \(u,q\in\mathbb R\) ist entsprechend
\(G_q(x)=q\RealMul x\) eine Selbstabbildung von \(\mathbb R\). Die Gleichungen
\[
  a_0=u,
  \qquad a_{\Succ(n)}=G_q(a_n)
\]
definieren daher eine eindeutig bestimmte reelle Folge, ohne dass eine
reelle Potenznotation vorweggenommen werden muss.
\end{example}

\chapter{Reelle Folgen und Grenzwerte}

\section{Reelle Notation und punktweise Operationen}

Band~15 konstruiert die reellen Operationen zunächst unter eigenen Symbolen.
In diesem und in späteren analytischen Bänden werden dafür die üblichen
Körpersymbole als reine Abkürzungen verwendet. Rationale Zahlen und
insbesondere rationale Zahlzeichen werden dabei, wie am Ende von Band~15
vereinbart, über \(\iota_{\mathbb R}\) mit ihren reellen Bildern
identifiziert.
Die folgenden Abkürzungen gelten innerhalb dieses Analysis-Kapitels, sobald
die beteiligten Terme als reell typisiert sind. Sie ändern nicht die
natürlichen Tags \(0,1,\ldots\) in späteren rein mengentheoretischen Kapiteln.

\FormulaDefDeltaK[Konventionelle Schreibweise für reelle Operationen]{%
  \begin{aligned}[t]
  &0\coloneqq\RealZero\dsep 1\coloneqq\RealOne,\\[-2pt]
  &x,y\in\mathbb R\vdash
    \begin{gathered}[t]
      x+y\coloneqq x\RealAdd y,\qquad
      -x\coloneqq\RealNeg{x},\qquad
      x-y\coloneqq x\RealAdd\RealNeg{y},\\[-2pt]
      xy\coloneqq x\RealMul y,\qquad
      (x\leq y)\coloneqq(x\RealLe y),\qquad
      (x<y)\coloneqq(x\RealLt y),\\[-2pt]
      (x\geq y)\coloneqq(y\RealLe x),\qquad
      (x>y)\coloneqq(y\RealLt x),\qquad
      |x|\coloneqq\RealAbs{x},
    \end{gathered}\\[-2pt]
  &x,y\in\mathbb R\dsep y\neq0\vdash
      y^{-1}\coloneqq\RealInv{y}\dsep
      \frac{x}{y}\coloneqq x\RealMul\RealInv{y}.
  \end{aligned}%
}{RealSequenceConventionalNotation}{%
  \DeltaRow{Mengen}{x\dsep y\dsep\mathbb R}
  \DeltaRow{Relationsoperatoren}{\RealLe\dsep\RealLt}
  \DeltaRow{Funktionsoperatoren}{\RealAdd\dsep\RealMul
    \dsep\RealNeg{\,\cdot\,}\dsep\RealInv{\,\cdot\,}
    \dsep\RealAbs{\,\cdot\,}}
}

Alle folgenden Rechnungen mit gewöhnlichen Symbolen sind daher Aussagen über
die in Band~15 registrierten Operationen; es wird keine zweite reelle
Arithmetik eingeführt.

\FormulaDefDeltaK[Punktweise Operationen reeller Folgen]{%
  \begin{aligned}[t]
  &a,b\colon\mathbb N\to\mathbb R\vdash
    a+b,\ a-b,\ a\mathbin{\cdot}b\colon\mathbb N\to\mathbb R\dsep{}\\[-2pt]
  &\quad\forall n\in\mathbb N\,
    \begin{gathered}[t]
      (a+b)_n\coloneqq a_n+b_n,\qquad
      (a-b)_n\coloneqq a_n-b_n,\\[-2pt]
      (a\mathbin{\cdot}b)_n\coloneqq a_nb_n
    \end{gathered}\\[-2pt]
  &a\colon\mathbb N\to\mathbb R\vdash
    -a,\ |a|\colon\mathbb N\to\mathbb R\dsep
    \forall n\in\mathbb N\,
      \bigl((-a)_n\coloneqq-a_n\dsep(|a|)_n\coloneqq|a_n|\bigr),\\[-2pt]
  &a\colon\mathbb N\to\mathbb R\dsep c\in\mathbb R\vdash
    ca\colon\mathbb N\to\mathbb R\dsep
    \forall n\in\mathbb N\,((ca)_n\coloneqq c\,a_n),\\[-2pt]
  &a\colon\mathbb N\to\mathbb R\dsep L\in\mathbb R\vdash
    a-L\colon\mathbb N\to\mathbb R\dsep
    \forall n\in\mathbb N\,((a-L)_n\coloneqq a_n-L).
  \end{aligned}%
}{PointwiseRealSequenceOperationsDef}{%
  \DeltaRow{Mengen}{c\dsep L\dsep n}
  \DeltaRow{Folgen}{a\dsep b\dsep a+b\dsep a-b}%
  \DeltaRow{Folgen}{-a\dsep ca\dsep a\mathbin{\cdot}b}%
  \DeltaRow{Folgen}{|a|\dsep a-L}
}
\begin{proof}
Die Abgeschlossenheitssätze
\FormulaRefAuto{RealCutAdditiveClosure}[thm] und
\FormulaRefAuto{RealCutMultiplicativeClosure}[thm] sowie
\FormulaRefAuto{RealAbsoluteValue}[def] zeigen für jeden Index einmal, dass
alle rechts stehenden Werte in \(\mathbb R\) liegen.
\FormulaRefAuto{IndexedFamilyTypedDefinitionPrinciple}[thm] liefert daraus
die angezeigten Funktionen; ihre Typisierung ist anschließend in den
Funktionsaussagen gespeichert.
\end{proof}

\FormulaDefDeltaK[Punktweiser Kehrwert und Quotient reeller Folgen]{%
  \begin{aligned}[t]
  &b\colon\mathbb N\to\mathbb R\dsep
    \forall n\in\mathbb N\,(b_n\neq0)\vdash\\[-2pt]
  &\quad b^{-1}\colon\mathbb N\to\mathbb R\dsep
    \forall n\in\mathbb N\,
      ((b^{-1})_n\coloneqq\tfrac1{b_n}),\\[-2pt]
  &a,b\colon\mathbb N\to\mathbb R\dsep
    \forall n\in\mathbb N\,(b_n\neq0)\vdash\\[-2pt]
  &\quad\frac ab\colon\mathbb N\to\mathbb R\dsep
    \forall n\in\mathbb N\,
      \left(\left(\frac ab\right)_n\coloneqq\frac{a_n}{b_n}\right).
  \end{aligned}%
}{PointwiseRealSequenceReciprocalQuotientDef}{%
  \DeltaRow{Mengen}{n}
  \DeltaRow{Funktionen}{a\dsep b\dsep b^{-1}\dsep a/b}
}
\begin{proof}
Die Nichtnullbedingung macht \(\RealInv{b_n}\) für jeden Index anwendbar.
\FormulaRefAuto{RealCutMultiplicativeClosure}[thm] und erneut das einmalige
Typisierungsprinzip liefern beide Funktionen.
\end{proof}

\section{Konvergenz}

Von nun an seien Folgen ohne anderslautenden Zusatz reelle Folgen, also
Funktionen \(a\colon\mathbb N\to\mathbb R\). Der Ausdruck
\(\varepsilon>0\) bedeutet stets
\(\varepsilon\in\mathbb R\land\RealZero\RealLt\varepsilon\).

\FormulaDefDeltaK[Konvergenz einer reellen Folge]{%
  \begin{aligned}[t]
  &a\colon\mathbb N\to\mathbb R\dsep L\in\mathbb R\vdash
    \bigl(a_n\longrightarrow L\bigr)\coloneqq{}\\[-2pt]
  &\quad\forall\varepsilon\in\mathbb R\,
    \Bigl(\varepsilon>0\rightarrow
      \exists N\in\mathbb N\,
      \forall n\in\mathbb N\,
        (N\leq n\rightarrow|a_n-L|<\varepsilon)\Bigr)
  \end{aligned}%
}{RealSequenceConvergenceDef}{%
  \DeltaRow{Mengen}{L\dsep\varepsilon\dsep N\dsep n}
  \DeltaRow{Funktionen}{a}
  \DeltaRow{Neue Schreibweise}{a_n\to L}
}

Gilt \(a_n\to L\), so heißt \(L\) ein \textbf{Grenzwert} der Folge. Man
schreibt auch
\[
  \lim_{n\to\infty}a_n=L.
\]
Die Zahl \(N\) darf von \(\varepsilon\) abhängen, aber nicht vom später
gewählten \(n\). Die Forderung betrifft nur Indizes \(n\geq N\); endlich
viele Anfangsglieder haben deshalb keinen Einfluss auf Konvergenz und
Grenzwert.

\FormulaDefDeltaK[Konvergente und divergente Folge]{%
  a\colon\mathbb N\to\mathbb R\vdash
  \operatorname{Konv}(a)\coloneqq
    \exists L\in\mathbb R\,(a_n\to L)\dsep
  \operatorname{Div}(a)\coloneqq\neg\operatorname{Konv}(a)%
}{RealSequenceConvergentDivergentDef}{%
  \DeltaRow{Mengen}{L}
  \DeltaRow{Funktionen}{a}
  \DeltaRow{Neue Prädikate}{\operatorname{Konv}\dsep\operatorname{Div}}
}

\FormulaThmDeltaK[Eindeutigkeit des Grenzwerts]{%
  a\colon\mathbb N\to\mathbb R\dsep
  L,K\in\mathbb R\dsep
  a_n\to L\dsep a_n\to K\vdash L=K%
}{RealSequenceLimitUnique}{%
  \DeltaRow{Mengen}{L\dsep K\dsep\varepsilon\dsep N\dsep n}
  \DeltaRow{Funktionen}{a}
}
\begin{proof}
Angenommen, \(L\neq K\), und setze
\(\varepsilon=|L-K|/3>0\). Für hinreichend großes \(n\) gelten zugleich
\(|a_n-L|<\varepsilon\) und \(|a_n-K|<\varepsilon\). Mit der
Dreiecksungleichung folgt
\[
  |L-K|\leq |L-a_n|+|a_n-K|
  <2\varepsilon=\frac23|L-K|,
\]
ein Widerspruch. Also ist \(L=K\).
\end{proof}

Wegen der Eindeutigkeit ist \(\lim_{n\to\infty}a_n\) ein eindeutig
bestimmter Ausdruck, sobald die Folge konvergiert.

\FormulaThmDeltaK[Konvergenz konstanter Folgen]{%
  c\in\mathbb R\vdash
  (\operatorname{const}_{\mathbb N,c})_n\to c%
}{ConstantRealSequenceConverges}{%
  \DeltaRow{Mengen}{c\dsep n}
  \DeltaRow{Funktionen}{\operatorname{const}_{\mathbb N,c}}
}

\FormulaThmDeltaK[Endliche Änderungen verändern den Grenzwert nicht]{%
  \begin{aligned}[t]
  &a,b\colon\mathbb N\to\mathbb R\dsep
    \exists N_0\in\mathbb N\,
      \forall n\in\mathbb N\,(N_0\leq n\rightarrow a_n=b_n)\\[-2pt]
  &\quad\vdash
    \forall L\in\mathbb R\,(a_n\to L\leftrightarrow b_n\to L)
  \end{aligned}%
}{FiniteModificationPreservesLimit}{%
  \DeltaRow{Mengen}{L\dsep N_0\dsep n}
  \DeltaRow{Funktionen}{a\dsep b}
}
\begin{proof}
Zu einer gegebenen Schranke \(N\) für eine der Folgen verwendet man für die
andere Folge den größeren der beiden Indizes \(N\) und \(N_0\).
\end{proof}

\FormulaThmDeltaK[Restfolgen haben denselben Grenzwert]{%
  a\colon\mathbb N\to\mathbb R\dsep r\in\mathbb N\dsep L\in\mathbb R
  \vdash
  a_n\to L\leftrightarrow a^{\langle r\rangle}_n\to L%
}{SequenceTailSameLimit}{%
  \DeltaRow{Mengen}{L\dsep r\dsep n}
  \DeltaRow{Funktionen}{a\dsep a^{\langle r\rangle}}
}

\FormulaThmDeltaK[Teilfolgen einer konvergenten Folge]{%
  \begin{aligned}[t]
  &a\colon\mathbb N\to\mathbb R\dsep L\in\mathbb R\dsep a_n\to L\dsep
    \varphi\colon\mathbb N\to\mathbb N\dsep{}\\[-2pt]
  &\quad\forall n\in\mathbb N\,
    \bigl(\varphi(n)<\varphi(\Succ(n))\bigr)
    \vdash a_{\varphi(n)}\to L
  \end{aligned}%
}{SubsequenceOfConvergentSequence}{%
  \DeltaRow{Mengen}{L\dsep n}
  \DeltaRow{Funktionen}{a\dsep\varphi}
}
\begin{proof}
Eine streng wachsende natürliche Indexfolge erfüllt \(n\leq\varphi(n)\).
Ist \(n\geq N\), so ist daher auch \(\varphi(n)\geq N\), und die
\(\varepsilon\)-Abschätzung der Ausgangsfolge kann bei diesem Index
angewandt werden.
\end{proof}

\section{Nullfolgen und Abschätzungen}

\FormulaDefDeltaK[Nullfolge]{%
  a\colon\mathbb N\to\mathbb R\vdash
  \bigl(a\text{ ist eine Nullfolge}\bigr)
  \coloneqq(a_n\to0)%
}{NullSequenceDef}{%
  \DeltaRow{Funktionen}{a}
}

\FormulaThmDeltaK[Konvergenz als Nullfolgeneigenschaft]{%
  a\colon\mathbb N\to\mathbb R\dsep L\in\mathbb R\vdash
  a_n\to L\leftrightarrow(a-L)_n\to0%
}{ConvergenceIffDifferenceNullSequence}{%
  \DeltaRow{Mengen}{L\dsep n}
  \DeltaRow{Funktionen}{a\dsep a-L}
}

\FormulaThmDeltaK[Majorantenkriterium für Nullfolgen]{%
  a,b\colon\mathbb N\to\mathbb R\dsep
  b_n\to0\dsep
  \exists N_0\in\mathbb N\,
    \forall n\in\mathbb N\,
      (N_0\leq n\rightarrow|a_n|\leq b_n)
  \vdash a_n\to0%
}{NullSequenceMajorantCriterion}{%
  \DeltaRow{Mengen}{N_0\dsep n}
  \DeltaRow{Funktionen}{a\dsep b}
}

\FormulaThmDeltaK[Einschließungsprinzip]{%
  \begin{aligned}[t]
  &a,b,c\colon\mathbb N\to\mathbb R\dsep L\in\mathbb R\dsep
    a_n\to L\dsep c_n\to L\dsep{}\\[-2pt]
  &\quad\exists N_0\in\mathbb N\,
    \forall n\in\mathbb N\,
      \bigl(N_0\leq n\rightarrow
        (a_n\leq b_n\land b_n\leq c_n)\bigr)
    \vdash b_n\to L
  \end{aligned}%
}{RealSequenceSqueezeTheorem}{%
  \DeltaRow{Mengen}{L\dsep N_0\dsep n}
  \DeltaRow{Funktionen}{a\dsep b\dsep c}
}
\begin{proof}
Für ein gegebenes \(\varepsilon>0\) liegen schließlich sowohl \(a_n\) als
auch \(c_n\) im Intervall
\((L-\varepsilon,L+\varepsilon)\). Ab dem Maximum dieser beiden Schranken
und \(N_0\) liegt wegen \(a_n\leq b_n\leq c_n\) auch \(b_n\) in diesem
Intervall.
\end{proof}

\section{Beschränktheit und Ordnung}

\FormulaDefDeltaK[Beschränkte reelle Folge]{%
  a\colon\mathbb N\to\mathbb R\vdash
  \operatorname{Beschr}(a)\coloneqq
  \exists C\in\mathbb R\,
    \bigl(C\geq0\land
      \forall n\in\mathbb N\,(|a_n|\leq C)\bigr)%
}{BoundedRealSequenceDef}{%
  \DeltaRow{Mengen}{C\dsep n}
  \DeltaRow{Funktionen}{a}
  \DeltaRow{Neues Prädikat}{\operatorname{Beschr}}
}

\FormulaThmDeltaK[Konvergente Folgen sind beschränkt]{%
  a\colon\mathbb N\to\mathbb R\dsep L\in\mathbb R\dsep a_n\to L
  \vdash\operatorname{Beschr}(a)%
}{ConvergentRealSequenceBounded}{%
  \DeltaRow{Mengen}{L\dsep C\dsep N\dsep n}
  \DeltaRow{Funktionen}{a}
}
\begin{proof}
Für \(\varepsilon=1\) gibt es ein \(N\), so dass
\(|a_n-L|<1\) für alle \(n\geq N\). Dann ist
\(|a_n|<|L|+1\) für alle diese Indizes. Die endlich vielen Werte mit
\(n<N\) besitzen ein Maximum ihrer Absolutbeträge. Das Maximum dieses
Wertes und \(|L|+1\) ist eine Schranke für die ganze Folge. Für \(N=0\)
entfällt der endliche Anfangsteil.
\end{proof}

\FormulaThmDeltaK[Ordnung bleibt im Grenzübergang erhalten]{%
  \begin{aligned}[t]
  &a,b\colon\mathbb N\to\mathbb R\dsep A,B\in\mathbb R\dsep
    a_n\to A\dsep b_n\to B\dsep{}\\[-2pt]
  &\quad\exists N_0\in\mathbb N\,
    \forall n\in\mathbb N\,(N_0\leq n\rightarrow a_n\leq b_n)
    \vdash A\leq B
  \end{aligned}%
}{RealSequenceLimitPreservesOrder}{%
  \DeltaRow{Mengen}{A\dsep B\dsep N_0\dsep n}
  \DeltaRow{Funktionen}{a\dsep b}
}
\begin{proof}
Wäre \(A>B\), so könnte man
\(\varepsilon=(A-B)/3>0\) wählen. Schließlich wären dann
\(a_n>A-\varepsilon>B+\varepsilon>b_n\), im Widerspruch zur angenommenen
Ungleichung.
\end{proof}

Aus einer schließlich strikten Ungleichung \(a_n<b_n\) folgt im Allgemeinen
nur \(A\leq B\), nicht \(A<B\). Beispielsweise gilt
\(0<1/(n+1)\), obwohl beide Seiten gegen \(0\) konvergieren.

\FormulaDefDeltaK[Monotone reelle Folgen]{%
  a\colon\mathbb N\to\mathbb R\vdash
  \begin{aligned}[t]
  \operatorname{Mon}_{\uparrow}(a)&\coloneqq
    \forall n\in\mathbb N\,(a_n\leq a_{\Succ(n)}),\\[-2pt]
  \operatorname{Mon}_{\downarrow}(a)&\coloneqq
    \forall n\in\mathbb N\,(a_{\Succ(n)}\leq a_n)
  \end{aligned}%
}{MonotoneRealSequenceDef}{%
  \DeltaRow{Mengen}{n}
  \DeltaRow{Funktionen}{a}
  \DeltaRow{Neue Prädikate}{\operatorname{Mon}_{\uparrow}\dsep
    \operatorname{Mon}_{\downarrow}}
}

\FormulaThmDeltaK[Monotone Konvergenz]{%
  \begin{aligned}[t]
  &a\colon\mathbb N\to\mathbb R\dsep
    \operatorname{Mon}_{\uparrow}(a)\dsep
    \exists C\in\mathbb R\,\forall n\in\mathbb N\,(a_n\leq C)
    \vdash
    a_n\to\sup_{\mathbb R,\RealLe}\bigl(a[\mathbb N]\bigr),\\[-2pt]
  &a\colon\mathbb N\to\mathbb R\dsep
    \operatorname{Mon}_{\downarrow}(a)\dsep
    \exists C\in\mathbb R\,\forall n\in\mathbb N\,(C\leq a_n)
    \vdash
    a_n\to\inf_{\mathbb R,\RealLe}\bigl(a[\mathbb N]\bigr).
  \end{aligned}%
}{MonotoneRealSequenceConvergence}{%
  \DeltaRow{Mengen}{C\dsep n\dsep a[\mathbb N]}
  \DeltaRow{Funktionen}{a}
  \DeltaRow{Relationsoperatoren}{\RealLe}
}
\begin{proof}
Sei zunächst \(a\) monoton wachsend und nach oben beschränkt. Wegen der
Vollständigkeit von \(\mathbb R\) existiert
\(s=\sup_{\mathbb R,\RealLe}(a[\mathbb N])\). Zu \(\varepsilon>0\) ist
\(s-\varepsilon\) keine obere Schranke der Wertemenge. Also gibt es ein
\(N\) mit \(s-\varepsilon<a_N\leq s\). Für \(n\geq N\) liefert die
Monotonie
\(s-\varepsilon<a_N\leq a_n\leq s<s+\varepsilon\). Somit
\(a_n\to s\). Im fallenden Fall existiert das Infimum
\(r=\inf_{\mathbb R,\RealLe}(a[\mathbb N])\) nach
\FormulaRefAuto{RealInfimumCharacterization}[thm]. Nun ist
\(r+\varepsilon\) keine untere Schranke. Daher gibt es ein \(N\) mit
\(r\leq a_N<r+\varepsilon\), und für \(n\geq N\) gilt
\(r\leq a_n\leq a_N<r+\varepsilon\). Also \(a_n\to r\).
\end{proof}

\section{Grenzwertregeln}

\FormulaThmDeltaK[Lineare Grenzwertregeln]{%
  a,b\colon\mathbb N\to\mathbb R\dsep
  A,B,c\in\mathbb R\dsep
  a_n\to A\dsep b_n\to B
  \vdash
  \begin{aligned}[t]
  (a+b)_n&\to A+B,\\[-2pt]
  (a-b)_n&\to A-B,\\[-2pt]
  (ca)_n&\to cA
  \end{aligned}%
}{RealSequenceLinearLimitRules}{%
  \DeltaRow{Mengen}{A\dsep B\dsep c\dsep n}
  \DeltaRow{Funktionen}{a\dsep b\dsep a+b\dsep a-b\dsep ca}
}
\begin{proof}
Für die Summe wählt man zu \(\varepsilon>0\) für beide Folgen die
Genauigkeit \(\varepsilon/2\) und verwendet
\[
 |(a_n+b_n)-(A+B)|
 \leq|a_n-A|+|b_n-B|.
\]
Die Differenz folgt durch Addition von \(-b_n\). Für \(c=0\) ist die
Skalarmultiplikation konstant null; für \(c\neq0\) verwendet man für
\(a_n\) die Genauigkeit \(\varepsilon/|c|\).
\end{proof}

\FormulaThmDeltaK[Produktregel für Grenzwerte]{%
  a,b\colon\mathbb N\to\mathbb R\dsep
  A,B\in\mathbb R\dsep
  a_n\to A\dsep b_n\to B
  \vdash (a\mathbin{\cdot}b)_n\to AB%
}{RealSequenceProductLimitRule}{%
  \DeltaRow{Mengen}{A\dsep B\dsep n}
  \DeltaRow{Funktionen}{a\dsep b\dsep a\mathbin{\cdot}b}
}
\begin{proof}
Die konvergente Folge \(a\) ist beschränkt; sei \(|a_n|\leq C\). Dann
\[
 |a_nb_n-AB|
 \leq |a_n|\,|b_n-B|+|B|\,|a_n-A|
 \leq C|b_n-B|+|B|\,|a_n-A|.
\]
Die beiden Summanden werden durch geeignete positive Genauigkeiten kleiner
als \(\varepsilon/2\). Falls \(C=0\) oder \(B=0\), entfällt der betreffende
Summand beziehungsweise wird eine beliebige positive Hilfsschranke benutzt.
\end{proof}

\FormulaThmDeltaK[Kehrwertregel für Grenzwerte]{%
  b\colon\mathbb N\to\mathbb R\dsep
  B\in\mathbb R\dsep b_n\to B\dsep B\neq0\dsep
  \forall n\in\mathbb N\,(b_n\neq0)
  \vdash (b^{-1})_n\to\frac1B%
}{RealSequenceReciprocalLimitRule}{%
  \DeltaRow{Mengen}{B\dsep n}
  \DeltaRow{Funktionen}{b\dsep b^{-1}}
}
\begin{proof}
Aus \(b_n\to B\) und \(B\neq0\) folgt schließlich
\(|b_n|>|B|/2\). Für diese Indizes gilt
\[
 \left|\frac1{b_n}-\frac1B\right|
 =\frac{|b_n-B|}{|b_n|\,|B|}
 <\frac{2}{|B|^2}|b_n-B|.
\]
Die Behauptung folgt aus der Konvergenz von \(b_n\) gegen \(B\).
\end{proof}

\FormulaThmDeltaK[Quotientenregel für Grenzwerte]{%
  \begin{aligned}[t]
  &a,b\colon\mathbb N\to\mathbb R\dsep A,B\in\mathbb R\dsep
    a_n\to A\dsep b_n\to B\dsep B\neq0\dsep{}\\[-2pt]
  &\quad\forall n\in\mathbb N\,(b_n\neq0)
    \vdash \left(\frac ab\right)_n\to\frac AB
  \end{aligned}%
}{RealSequenceQuotientLimitRule}{%
  \DeltaRow{Mengen}{A\dsep B\dsep n}
  \DeltaRow{Funktionen}{a\dsep b\dsep a/b}
}
\begin{proof}
Nach den punktweisen Definitionen stimmen \(a/b\) und
\(a\mathbin{\cdot}b^{-1}\) an jedem Index überein und sind daher dieselbe
Funktion. Man wendet
\FormulaRefAuto{RealSequenceReciprocalLimitRule}[thm] auf \(b\) und
anschließend \FormulaRefAuto{RealSequenceProductLimitRule}[thm] an.
\end{proof}

\FormulaThmDeltaK[Stetigkeit des Betrags für Folgen]{%
  a\colon\mathbb N\to\mathbb R\dsep A\in\mathbb R\dsep a_n\to A
  \vdash (|a|)_n\to|A|%
}{AbsoluteValuePreservesSequenceLimits}{%
  \DeltaRow{Mengen}{A\dsep n}
  \DeltaRow{Funktionen}{a\dsep |a|}
}
\begin{proof}
Die umgekehrte Dreiecksungleichung liefert
\(\bigl||a_n|-|A|\bigr|\leq|a_n-A|\).
\end{proof}

\section{Cauchy-Folgen}

\FormulaDefDeltaK[Cauchy-Folge in den reellen Zahlen]{%
  \begin{aligned}[t]
  &a\colon\mathbb N\to\mathbb R\vdash
    \operatorname{Cauchy}(a)\coloneqq{}\\[-2pt]
  &\quad\forall\varepsilon\in\mathbb R\,
    \Bigl(\varepsilon>0\rightarrow
      \exists N\in\mathbb N\,
      \forall m,n\in\mathbb N\,\bigl((N\leq m\land N\leq n)\\[-2pt]
  &\hspace{54mm}\rightarrow|a_m-a_n|<\varepsilon\bigr)\Bigr)
  \end{aligned}%
}{RealCauchySequenceDef}{%
  \DeltaRow{Mengen}{\varepsilon\dsep N\dsep m\dsep n}
  \DeltaRow{Funktionen}{a}
  \DeltaRow{Neues Prädikat}{\operatorname{Cauchy}}
}

Die Cauchy-Bedingung erwähnt keinen vermuteten Grenzwert. Sie verlangt, dass
späte Folgenglieder untereinander beliebig nahe liegen. Ob daraus ein
Grenzwert folgt, hängt von der Vollständigkeit der Zielmenge ab.

\FormulaThmDeltaK[Konvergenz impliziert die Cauchy-Eigenschaft]{%
  a\colon\mathbb N\to\mathbb R\dsep L\in\mathbb R\dsep a_n\to L
  \vdash\operatorname{Cauchy}(a)%
}{ConvergentSequenceIsCauchy}{%
  \DeltaRow{Mengen}{L\dsep\varepsilon\dsep N\dsep m\dsep n}
  \DeltaRow{Funktionen}{a}
}
\begin{proof}
Zu \(\varepsilon>0\) wähle \(N\) so, dass
\(|a_k-L|<\varepsilon/2\) für alle \(k\geq N\). Für \(m,n\geq N\) gilt
dann
\[
 |a_m-a_n|\leq|a_m-L|+|a_n-L|<\varepsilon.
\]
\end{proof}

\FormulaThmDeltaK[Cauchy-Folgen sind beschränkt]{%
  a\colon\mathbb N\to\mathbb R\dsep\operatorname{Cauchy}(a)
  \vdash\operatorname{Beschr}(a)%
}{RealCauchySequenceBounded}{%
  \DeltaRow{Mengen}{N\dsep n}
  \DeltaRow{Funktionen}{a}
}
\begin{proof}
Wende die Cauchy-Bedingung mit \(\varepsilon=1\) an und fixiere einen Index
\(N\), ab dem sie gilt. Dann ist \(|a_n-a_N|<1\) für \(n\geq N\), also
\(|a_n|<|a_N|+1\). Die endlich vielen Anfangsglieder werden wie im Beweis
von \FormulaRefAuto{ConvergentRealSequenceBounded}[thm] gemeinsam
beschränkt.
\end{proof}

\FormulaThmDeltaK[Cauchy-Kriterium in den reellen Zahlen]{%
  a\colon\mathbb N\to\mathbb R\vdash
  \operatorname{Konv}(a)\leftrightarrow\operatorname{Cauchy}(a)%
}{RealSequenceCauchyCriterion}{%
  \DeltaRow{Funktionen}{a}
}
\begin{proof}
Die Hinrichtung ist
\FormulaRefAuto{ConvergentSequenceIsCauchy}[thm]. Sei umgekehrt \(a\) eine
Cauchy-Folge. Nach
\FormulaRefAuto{RealCauchySequenceBounded}[thm] ist sie beschränkt. Für
\(N\in\mathbb N\) setze
\[
  T_N\coloneqq
  \{a_n\mid n\in\mathbb N\land N\leq n\}.
\]
Jede Menge \(T_N\) ist nichtleer und nach unten beschränkt. Nach
\FormulaRefAuto{RealInfimumCharacterization}[thm] besitzt sie daher ein
eindeutiges Infimum
\(\alpha_N=\inf_{\mathbb R,\RealLe}(T_N)\in\mathbb R\). Das einmalige
Typisierungsprinzip liefert somit eine Folge
\(\alpha\colon\mathbb N\to\mathbb R\). Wegen
\(T_{\Succ(N)}\subseteq T_N\) ist \(\alpha_N\leq\alpha_{\Succ(N)}\).
Außerdem ist die Menge
\[
  \mathcal A\coloneqq\{\alpha_N\mid N\in\mathbb N\}
\]
nichtleer und durch jede obere Schranke von \(a[\mathbb N]\) nach oben
beschränkt. Ihre reelle Supremumszahl
\[
  L\coloneqq\sup_{\mathbb R,\RealLe}(\mathcal A)
\]
existiert folglich nach
\FormulaRefAuto{RealLeastUpperBoundProperty}[thm].

Sei nun \(\varepsilon>0\) und
\(\delta=\varepsilon/2\). Wähle \(N\) so, dass
\(|a_m-a_n|<\delta\) für alle \(m,n\geq N\). Für festes \(n\geq N\)
ist \(a_n-\delta\) eine untere Schranke von \(T_N\); daher
\[
  a_n-\delta\leq\alpha_N\leq L.
\]
Für jedes \(k\in\mathbb N\) kann man einen Index \(m\geq k,N\) wählen.
Dann gilt
\(\alpha_k\leq a_m<a_n+\delta\). Somit ist \(a_n+\delta\) eine obere
Schranke von \(\mathcal A\), also
\[
  L\leq a_n+\delta.
\]
Zusammen folgt \(|a_n-L|\leq\delta<\varepsilon\) für alle \(n\geq N\).
Damit konvergiert \(a\) gegen \(L\).
\end{proof}

Die Rückrichtung des Cauchy-Kriteriums ist genau die Stelle, an der die
Vollständigkeit von \(\mathbb R\) benötigt wird. In \(\mathbb Q\) gibt es
Cauchy-Folgen, deren reeller Grenzwert irrational ist und die deshalb in
\(\mathbb Q\) keinen Grenzwert besitzen.

\chapter{Die Mogiljanskaja-Familien}

\section{Die beiden Grundfamilien}

Für die spätere Konstruktion zweier unendlicher Halbgruppen werden zwei
disjunkte getaggte Schichten benötigt. In diesem Band geht es ausschließlich
um ihre familien- und mengentheoretischen Eigenschaften; Operationen auf den
späteren Trägermengen werden erst im Halbgruppenband definiert.
Die Zeichen \(0\) und \(1\) in den folgenden geordneten Paaren sind dabei
ausdrücklich die natürlichen Peano-Zahlen und dienen lediglich als Tags.

\FormulaDefDeltaK[Die einfach indizierte Grundschicht]{%
  L\coloneqq\{0\}\times\mathbb N%
}{MogiljanskajaLayerLDef}{%
  \DeltaRow{Mengen}{L}
}

\FormulaDefDeltaK[Die zweifach indizierte Grundschicht]{%
  D\coloneqq\{1\}\times(\mathbb N\times\mathbb N)%
}{MogiljanskajaLayerDDef}{%
  \DeltaRow{Mengen}{D}
}

\FormulaDefDeltaK[Die einfach indizierte Grundfamilie]{%
  a\colon\mathbb N\to L\dsep
  i\in\mathbb N\vdash a_i\coloneqq(0,i)%
}{MogiljanskajaAElementsDef}{%
  \DeltaRow{Mengen}{L\dsep a_i\dsep i}
  \DeltaRow{Funktionen}{a}
}

Die Typisierung ist unmittelbar, denn aus \(i\in\mathbb N\) folgt
\((0,i)\in\{0\}\times\mathbb N=L\). Nach dem einmaligen
Typisierungsprinzip ist damit die ganze Folge definiert.

\FormulaDefDeltaK[Die zweifach indizierte Grundfamilie]{%
  b\colon\mathbb N\times\mathbb N\to D\dsep
  i,j\in\mathbb N\vdash b_{ij}\coloneqq(1,(i,j))%
}{MogiljanskajaBElementsDef}{%
  \DeltaRow{Mengen}{D\dsep b_{ij}\dsep i\dsep j}
  \DeltaRow{Funktionen}{b}
}

Hier steht \(b_{ij}\) für \(b((i,j))\). Die Familie ist also eine einzige
Funktion auf \(\mathbb N\times\mathbb N\), keine zweistufige Sammlung
voneinander unabhängiger Funktionen.

\FormulaThmDeltaK[Eindeutigkeit der (a)-Indizes]{%
  i,j\in\mathbb N\dsep a_i=a_j\vdash i=j%
}{MogiljanskajaAIndexUnique}{%
  \DeltaRow{Mengen}{a_i\dsep a_j\dsep i\dsep j}
  \DeltaRow{Funktionen}{a}
}
\begin{proof}
Aus \((0,i)=(0,j)\) folgt durch Injektivität des geordneten Paars
\(i=j\).
\end{proof}

\FormulaThmDeltaK[Eindeutigkeit der (b)-Indizes]{%
  i,j,m,n\in\mathbb N\dsep b_{ij}=b_{mn}
  \vdash i=m\land j=n%
}{MogiljanskajaBIndexUnique}{%
  \DeltaRow{Mengen}{b_{ij}\dsep b_{mn}\dsep i\dsep j\dsep m\dsep n}
  \DeltaRow{Funktionen}{b}
}
\begin{proof}
Aus \((1,(i,j))=(1,(m,n))\) folgt zunächst \((i,j)=(m,n)\) und daraus
\(i=m\) sowie \(j=n\).
\end{proof}

\FormulaThmDeltaK[Elementkriterium der (a)-Schicht]{%
  s\in L\leftrightarrow\exists i\in\mathbb N\,(s=a_i)%
}{MogiljanskajaLayerLMembership}{%
  \DeltaRow{Mengen}{L\dsep s\dsep a_i\dsep i}
  \DeltaRow{Funktionen}{a}
}
\begin{proof}
Jedes Element von \(L=\{0\}\times\mathbb N\) besitzt eindeutig die Form
\((0,i)=a_i\), und jedes \(a_i\) liegt nach seiner Definition in \(L\).
\end{proof}

\FormulaThmDeltaK[Elementkriterium der (b)-Schicht]{%
  s\in D\leftrightarrow
    \exists i,j\in\mathbb N\,(s=b_{ij})%
}{MogiljanskajaLayerDMembership}{%
  \DeltaRow{Mengen}{D\dsep s\dsep b_{ij}\dsep i\dsep j}
  \DeltaRow{Funktionen}{b}
}
\begin{proof}
Jedes Element von
\(D=\{1\}\times(\mathbb N\times\mathbb N)\) besitzt eindeutig die Form
\((1,(i,j))=b_{ij}\), und jedes \(b_{ij}\) liegt in \(D\).
\end{proof}

\FormulaThmDeltaK[Bildmenge der (a)-Familie]{%
  a[\mathbb N]=L%
}{MogiljanskajaLayerLImage}{%
  \DeltaRow{Mengen}{L\dsep a[\mathbb N]}
  \DeltaRow{Funktionen}{a}
}
\begin{proof}
Die beiden Mengen besitzen nach
\FormulaRefAuto{MogiljanskajaLayerLMembership}[thm] dieselben Elemente.
\end{proof}

\FormulaThmDeltaK[Bildmenge der (b)-Familie]{%
  b[\mathbb N\times\mathbb N]=D%
}{MogiljanskajaLayerDImage}{%
  \DeltaRow{Mengen}{D\dsep b[\mathbb N\times\mathbb N]}
  \DeltaRow{Funktionen}{b}
}
\begin{proof}
Die beiden Mengen besitzen nach
\FormulaRefAuto{MogiljanskajaLayerDMembership}[thm] dieselben Elemente.
\end{proof}

\FormulaThmDeltaK[Die beiden Grundschichten sind disjunkt]{%
  L\cap D=\varnothing%
}{MogiljanskajaLayersDisjoint}{%
  \DeltaRow{Mengen}{L\dsep D}
}
\begin{proof}
Ein Element von \(L\) hat den ersten Tag \(0\), ein Element von \(D\) den
ersten Tag \(1\). Wegen \(0\neq1\) kann kein Element in beiden Schichten
liegen.
\end{proof}

\section{Marker und Reserveschichten}

\FormulaDefDeltaK[Marker und unendliche Reserveschicht]{%
  \begin{aligned}[t]
    c_0&\coloneqq b_{00},&
    c_1&\coloneqq b_{11},&
    c_2&\coloneqq b_{10},\\[-2pt]
    U&\coloneqq\bigl\{u\in D\bigm|
      \exists n,j\in\mathbb N\,
        u=b_{\Succ(\Succ(n)),j}\bigr\},\\[-2pt]
    V&\coloneqq U\cup\{c_2\},&
    P&\coloneqq\{c_0,c_1\},&
    D'&\coloneqq P\cup V
  \end{aligned}%
}{MogiljanskajaDPrimeDef}{%
  \DeltaRow{Mengen}{D\dsep D'\dsep U\dsep V\dsep P\dsep
    c_0\dsep c_1\dsep c_2\dsep b_{ij}\dsep n\dsep j\dsep u}
  \DeltaRow{Funktionen}{b}
}

Somit besteht \(U\) genau aus den Elementen der Zeilen mit erstem Index
mindestens \(2\). Die Menge \(V\) adjungiert den Marker \(c_2=b_{10}\),
und \(D'\) adjungiert zusätzlich \(c_0=b_{00}\) und \(c_1=b_{11}\).

\FormulaThmDeltaK[Elementkriterium der Reserveschicht]{%
  u\in U\eqvdash
  \exists n,j\in\mathbb N\,
    (u=b_{\Succ(\Succ(n)),j})%
}{MogiljanskajaReserveMembership}{%
  \DeltaRow{Mengen}{D\dsep U\dsep u\dsep b_{ij}\dsep n\dsep j}
  \DeltaRow{Funktionen}{b}
}

\FormulaThmDeltaK[Die Marker und das ausgesparte Rechteckelement]{%
  \begin{aligned}[t]
    &c_0,c_1,c_2\in D\dsep
      c_0,c_1,c_2\notin U\dsep
      c_0,c_1,c_2\text{ paarweise verschieden},\\[-2pt]
    &D'\subseteq D\dsep b_{01}\notin D'
  \end{aligned}%
}{MogiljanskajaReserveMarkerFacts}{%
  \DeltaRow{Mengen}{D\dsep D'\dsep U\dsep V\dsep P\dsep
    c_0\dsep c_1\dsep c_2\dsep b_{01}\dsep n\dsep j}
  \DeltaRow{Funktionen}{b}
}
\begin{proof}
Die drei Marker besitzen die verschiedenen Indexpaare
\[
  (0,0),\qquad(1,1),\qquad(1,0).
\]
Ihre Verschiedenheit folgt aus
\FormulaRefAuto{MogiljanskajaBIndexUnique}[thm]. Keiner hat einen ersten
Index von mindestens \(2\), also liegen sie außerhalb von \(U\). Folglich
liegen \(U,V,P,D'\) in \(D\). Das Element \(b_{01}\) gehört weder zu
\(P\) noch zu \(V\): Es ist keiner der Marker und sein erster Index ist
\(0\), so dass es nicht in \(U\) liegt.
\end{proof}

\FormulaThmDeltaK[Der dritte Marker liegt außerhalb des festen Kerns]{%
  c_2\notin P\dsep P\cap U=\varnothing%
}{MogiljanskajaReserveCoreFacts}{%
  \DeltaRow{Mengen}{P\dsep U\dsep c_0\dsep c_1\dsep c_2}
}

\section{Parametrisierungen und ihre Bildmengen}

\FormulaDefDeltaK[Die beiden Parametrisierungen der Produktschichten]{%
  \begin{aligned}[t]
    \delta&\colon\mathbb N\times\mathbb N\to D,&
      \delta(n,j)&\coloneqq b_{nj},\\[-2pt]
    \upsilon&\colon\mathbb N\times\mathbb N\to U,&
      \upsilon(n,j)&\coloneqq b_{\Succ(\Succ(n)),j}
  \end{aligned}%
}{MogiljanskajaRowParametrizationsDef}{%
  \DeltaRow{Mengen}{D\dsep U\dsep b_{ij}\dsep n\dsep j}
  \DeltaRow{Funktionen}{b\dsep\delta\dsep\upsilon}
}

\FormulaThmDeltaK[Die Zeilenparametrisierungen sind bijektiv]{%
  \delta\colon\mathbb N\times\mathbb N\bij D\dsep
  \upsilon\colon\mathbb N\times\mathbb N\bij U%
}{MogiljanskajaRowParametrizationsBijective}{%
  \DeltaRow{Mengen}{D\dsep U\dsep n\dsep j\dsep m\dsep q\dsep u}
  \DeltaRow{Funktionen}{\delta\dsep\upsilon}
}
\begin{proof}
Die Bijektivität von \(\delta\) ist die eindeutige Darstellung jedes
Elements von \(D\) als \(b_{nj}\). Für \(\upsilon\) liefert
\FormulaRefAuto{MogiljanskajaReserveMembership}[thm] die Surjektivität.
Aus
\(b_{\Succ(\Succ(n)),j}=b_{\Succ(\Succ(m)),q}\) folgen mit der
Indexeindeutigkeit zunächst
\(\Succ(\Succ(n))=\Succ(\Succ(m))\) und \(j=q\), dann durch zweimalige
Nachfolgerinjektivität \(n=m\).
\end{proof}

Insbesondere gelten die präzisen Bildmengengleichheiten
\[
  \delta[\mathbb N\times\mathbb N]=D,
  \qquad
  \upsilon[\mathbb N\times\mathbb N]=U.
\]

\FormulaDefDeltaK[Die Zeilenrückverschiebung]{%
  \theta\coloneqq\delta\circ\upsilon^{-1}%
}{MogiljanskajaThetaDef}{%
  \DeltaRow{Mengen}{D\dsep U}
  \DeltaRow{Funktionen}{\delta\dsep\upsilon\dsep\theta}
}

\FormulaThmDeltaK[Die Zeilenrückverschiebung ist bijektiv]{%
  \theta\colon U\bij D%
}{MogiljanskajaThetaBijective}{%
  \DeltaRow{Mengen}{D\dsep U}
  \DeltaRow{Funktionen}{\delta\dsep\upsilon\dsep\theta}
}
\begin{proof}
Inverse und Komposition bijektiver Funktionen sind bijektiv.
\end{proof}

\section{Die punktierte Verschiebung}

\FormulaThmDeltaK[Die Werte der Verschiebungsfolge liegen in der Reserve]{%
  n\in\mathbb N\vdash b_{\Succ(n),0}\in V%
}{MogiljanskajaShiftValueInV}{%
  \DeltaRow{Mengen}{U\dsep V\dsep c_2\dsep b_{ij}\dsep n}
  \DeltaRow{Funktionen}{b}
}
\begin{proof}
Für \(n=0\) ist \(b_{\Succ(0),0}=b_{10}=c_2\in V\). Für \(n>0\)
existiert ein \(m\in\mathbb N\) mit \(n=\Succ(m)\); dann ist
\(b_{\Succ(n),0}=b_{\Succ(\Succ(m)),0}\in U\subseteq V\).
\end{proof}

\FormulaDefDeltaK[Die Folge für die punktierte Verschiebung]{%
  H\colon\mathbb N\to V\dsep
  n\in\mathbb N\vdash H_n\coloneqq b_{\Succ(n),0}%
}{MogiljanskajaShiftSequenceDef}{%
  \DeltaRow{Mengen}{V\dsep b_{ij}\dsep n}
  \DeltaRow{Funktionen}{b\dsep H}
}

\FormulaThmDeltaK[Die punktierte Verschiebung entfernt genau den dritten
Marker]{%
  H\colon\mathbb N\inj V\dsep H_0=c_2\dsep
  V\setminus\{c_2\}=U%
}{MogiljanskajaShiftSequenceFacts}{%
  \DeltaRow{Mengen}{U\dsep V\dsep c_2\dsep n\dsep m}
  \DeltaRow{Funktionen}{H}
}
\begin{proof}
Aus \(H_n=H_m\) folgt
\(b_{\Succ(n),0}=b_{\Succ(m),0}\), also \(\Succ(n)=\Succ(m)\) und
damit \(n=m\). Ferner ist \(H_0=b_{10}=c_2\). Da \(c_2\notin U\) und
\(V=U\cup\{c_2\}\), gilt \(V\setminus\{c_2\}=U\).
\end{proof}

\FormulaThmDeltaK[Tatsächliche Bildmenge der Verschiebungsfolge]{%
  H[\mathbb N]
  =\{b_{\Succ(n),0}\mid n\in\mathbb N\}
  \dsep H[\mathbb N]\subsetneq V%
}{MogiljanskajaShiftSequenceImage}{%
  \DeltaRow{Mengen}{V\dsep b_{ij}\dsep n}
  \DeltaRow{Funktionen}{H}
}
\begin{proof}
Die erste Gleichheit ist das Bildmengenkriterium für die definierende
Vorschrift von \(H\). Die Inklusion in \(V\) folgt aus der Typisierung.
Sie ist echt: Das Element \(b_{\Succ(\Succ(0)),1}=b_{21}\) liegt in
\(U\subseteq V\). Jedes Glied \(H_n=b_{\Succ(n),0}\) besitzt dagegen den
zweiten Index \(0\); wegen der Eindeutigkeit der \(b\)-Indizes kann
\(b_{21}\) kein Glied von \(H\) sein.
\end{proof}

Damit ist ausdrücklich zwischen Ziel- und Bildmenge unterschieden:
\(H\colon\mathbb N\to V\), aber \(H[\mathbb N]\neq V\). Die Folge wird
nicht zur Aufzählung von \(V\), sondern zur Konstruktion einer Verschiebung
auf \(V\) verwendet.

\FormulaDefDeltaK[Die Bijektion in die um einen Marker erweiterte Reserve]{%
  \sigma\coloneqq(\SuccShift{H})^{-1}%
}{MogiljanskajaSigmaDef}{%
  \DeltaRow{Mengen}{U\dsep V}
  \DeltaRow{Funktionen}{H\dsep\SuccShift{H}\dsep\sigma}
}

\FormulaThmDeltaK[Die Markeradjunktion ist bijektiv]{%
  \sigma\colon U\bij V%
}{MogiljanskajaSigmaBijective}{%
  \DeltaRow{Mengen}{U\dsep V\dsep c_2}
  \DeltaRow{Funktionen}{H\dsep\SuccShift{H}\dsep\sigma}
}
\begin{proof}
Nach \FormulaRefAuto{MogiljanskajaShiftSequenceFacts}[thm] ist
\(H\colon\mathbb N\inj V\), \(H_0=c_2\) und
\(V\setminus\{c_2\}=U\). Der allgemeine Satz über die Folgenverschiebung
entlang einer Einbettung liefert
\(\SuccShift{H}\colon V\bij V\setminus\{H_0\}=U\). Daher ist die inverse
Funktion \(\sigma\colon U\bij V\).
\end{proof}

\section{Familien mit den vorgegebenen Bildmengen}

Die Familien \(a\), \(b\), \(\delta\) und \(\upsilon\) parametrisieren
bereits \(L,D,D\) beziehungsweise \(U\). Aus \(\sigma\) entsteht ohne eine
weitere Fallunterscheidung eine bijektive Familie mit Bild \(V\).

\FormulaDefDeltaK[Eine bijektive Familie mit Bildmenge (V)]{%
  \nu\coloneqq\sigma\circ\upsilon%
}{MogiljanskajaVParametrizationDef}{%
  \DeltaRow{Mengen}{U\dsep V}
  \DeltaRow{Funktionen}{\upsilon\dsep\sigma\dsep\nu}
}

\FormulaThmDeltaK[Die Familie (nu) parametrisiert (V)]{%
  \nu\colon\mathbb N\times\mathbb N\bij V\dsep
  \nu[\mathbb N\times\mathbb N]=V%
}{MogiljanskajaVParametrizationBijective}{%
  \DeltaRow{Mengen}{V}
  \DeltaRow{Funktionen}{\nu}
}
\begin{proof}
Die Funktionen \(\upsilon\colon\mathbb N^2\bij U\) und
\(\sigma\colon U\bij V\) sind bijektiv, also auch ihre Komposition.
\end{proof}

Für \(D'=P\cup V\) müssen die beiden Marker aus \(P\) und die Familie
\(\nu\) disjunkt zusammengefügt werden. Die Tags \(0\) und \(1\) sorgen
dafür, dass sich die beiden Indexteile nicht überschneiden.

\FormulaDefDeltaK[Indexmenge für die Familie mit Bild (D')]{%
  I_{D'}\coloneqq
  (\{0\}\times\NatSeg{2})
  \cup
  (\{1\}\times(\mathbb N\times\mathbb N))%
}{MogiljanskajaDPrimeIndexSetDef}{%
  \DeltaRow{Mengen}{I_{D'}}
}

\FormulaDefDeltaK[Eine Familie mit Bildmenge (D')]{%
  \begin{aligned}[t]
  \omega&\colon I_{D'}\to D',\\[-2pt]
  \omega((0,0))&\coloneqq c_0,\\[-2pt]
  \omega((0,1))&\coloneqq c_1,\\[-2pt]
  n,j\in\mathbb N\vdash
  \omega((1,(n,j)))&\coloneqq\nu(n,j)
  \end{aligned}%
}{MogiljanskajaDPrimeParametrizationDef}{%
  \DeltaRow{Mengen}{I_{D'}\dsep D'\dsep c_0\dsep c_1\dsep n\dsep j}
  \DeltaRow{Funktionen}{\nu\dsep\omega}
}

\FormulaThmDeltaK[Die Familie (omega) parametrisiert (D')]{%
  \omega\colon I_{D'}\bij D'\dsep \omega[I_{D'}]=D'%
}{MogiljanskajaDPrimeParametrizationBijective}{%
  \DeltaRow{Mengen}{I_{D'}\dsep D'\dsep P\dsep V}
  \DeltaRow{Funktionen}{\omega}
}
\begin{proof}
Auf dem ersten Indexteil nimmt \(\omega\) genau die beiden verschiedenen
Werte \(c_0,c_1\) aus \(P\) an. Auf dem zweiten Indexteil ist \(\omega\)
die bijektive Familie \(\nu\) mit Bild \(V\). Nach
\FormulaRefAuto{MogiljanskajaReserveMarkerFacts}[thm] und
\FormulaRefAuto{MogiljanskajaReserveCoreFacts}[thm] gilt
\(P\cap V=\varnothing\). Daher sind die beiden Wertebereiche disjunkt, und
die zusammengefügte Familie ist bijektiv auf \(P\cup V=D'\).
\end{proof}

\FormulaThmDeltaK[Mogiljanskaja-Bildmengen]{%
  \begin{aligned}[t]
    a[\mathbb N]&=L,&
    b[\mathbb N\times\mathbb N]&=D,\\[-2pt]
    \upsilon[\mathbb N\times\mathbb N]&=U,&
    \nu[\mathbb N\times\mathbb N]&=V,\\[-2pt]
    \omega[I_{D'}]&=D'
  \end{aligned}%
}{MogiljanskajaFamilyImages}{%
  \DeltaRow{Mengen}{L\dsep D\dsep U\dsep V\dsep D'\dsep I_{D'}}
  \DeltaRow{Funktionen}{a\dsep b\dsep\upsilon\dsep\nu\dsep\omega}
}

Diese Zusammenfassung ist die Schnittstelle zum späteren Halbgruppenband.
Dort müssen die Familien und ihre Typisierung nicht erneut aufgebaut werden;
es genügt, ihre Bildmengengleichheiten und die wenigen benötigten
Markerfakten zu zitieren.

\end{document}
